% ============================================================
%  CONCLUSION GÉNÉRALE
% ============================================================
\chapter*{Conclusion générale et perspectives}
\addcontentsline{toc}{chapter}{Conclusion générale et perspectives}
\markboth{Conclusion générale}{Conclusion générale}

\hspace{1cm}Ce projet de fin d'études a porté sur \textit{[résumer le sujet du projet]}. Au terme de ce travail, nous pouvons dresser le bilan suivant.

\vspace{0.5cm}
\textbf{Bilan des réalisations :}

\begin{itemize}
  \item Une étude approfondie des techniques d'\gls{ia} appliquées à la cybersécurité a été menée, permettant d'identifier les approches les plus prometteuses.
  \item Un modèle de détection basé sur les \gls{cnn} a été conçu, implémenté et évalué.
  \item Les résultats obtenus montrent une précision de 96,3\%, surpassant les méthodes de l'état de l'art.
\end{itemize}

\vspace{0.5cm}
\textbf{Perspectives :}

\begin{enumerate}
  \item Intégrer des mécanismes d'apprentissage par transfert pour améliorer la généralisation du modèle.
  \item Explorer les architectures de type Transformer pour le traitement séquentiel du trafic réseau.
  \item Déployer le système dans un environnement de production avec une infrastructure \gls{soc} complète.
  \item Étendre la solution à la détection de menaces dans les environnements \gls{iot}.
\end{enumerate}
